% Part: first-order-logic
% Chapter: beyond
% Section: higher-order-logic

\documentclass[../../../include/open-logic-section]{subfiles}

\begin{document}

\olfileid{fol}{byd}{hol}

\olsection{اعلیٰ درجے دی منطق}

پہلے درجے دی منطق توں دوجے درجے دی منطق
ول آ کے رسمی زبان دے اندر پہلے درجے دے
!!{domain} دیاں شےواں دے سیٹاں بارے
گل کرنا ممکن ہویا۔ فیر ایتھے ہی کیوں
رکیے؟ مثال لئی تیجے درجے دی منطق نال
شےواں دے سیٹاں دے سیٹ وچارے جا سکدے
نیں، یا شاید ایسے سیٹ وی جیہناں وچ
شےواں تے شےواں دے سیٹ دوویں شامل ہون۔
چوتھے درجے دی منطق فیر ایس قسم دیاں
شےواں دے سیٹاں بارے گل کرن دی
اجازت دیوے گی۔ جداں تسی شاید
اندازہ لا لیا اے، ایہ خیال
من چاہے درجے تک دہرایا جا سکدا اے۔

عمل وچ اعلیٰ درجے دی منطق دی
!!{formula} بندی اکثر تعلقاں دی بجائے فنکشناں
دے ذریعے کیتی جاندی اے۔
(قدرتی مطابقتاں نوں ملحوظ رکھ کے
ایہ فرق بنیادی نہیں۔) کجھ بنیادی
”قسماں“ $A$, $B$, $C$,~\dots
دتی ہوون (جیہناں نوں ہن ”انواع“
یا types آکھاں گے)، تاں اسیں
ہیٹھلی شرط نال نویاں انواع بنا
سکدے آں:
\begin{quote}
جے $\sigma$ تے $\tau$ محدود
تشکیل والیاں انواع ہون،
تاں $\sigma \to \tau$ وی ایسی
نوع اے۔
\end{quote}
انواع نوں نحوی ”نشان“ سمجھو
جیہڑے ساڈے !!{domain} دیاں
شےواں دی قسم بندی کردے نیں؛
$\sigma \to \tau$ انہاں شےواں
نوں بیان کردا اے جیہڑیاں نوع~$\sigma$
دیاں شےواں توں نوع~$\tau$ دیاں
شےواں ول فنکشن نیں۔ مثال لئی
سچائی دیاں قدراں ”سچ“ تے ”جھوٹ“
دی اک نوع $\Omega$، تے قدرتی
اعداد دی اک نوع $\Nat$ رکھ
سکدے آں۔ فیر نوع $\Nat \to
\Omega$ دیاں شےواں نوں اک-جگہی
تعلق یا $\Nat$ دے ذیلی سیٹ
سمجھیا جا سکدا اے؛ نوع
$\Nat \to \Nat$ دیاں شےواں
قدرتی اعداد توں قدرتی اعداد
ول فنکشن نیں؛ تے نوع
$(\Nat \to \Nat) \to \Nat$
دیاں شےواں ”فنکشنل“
نیں، یعنی اعلیٰ نوع دے
اوہ فنکشن جیہڑے فنکشناں نوں
اعداد ول بھیج دے نیں۔

دوجے درجے دی منطق وانگ،
اعلیٰ درجے دی منطق نوں
کئی قسماں والی منطق دا
اک روپ سمجھیا جا سکدا اے،
جتھے ہر مطلوبہ نوع دیاں
شےواں لئی اک قسم ہوئے۔
پر عموماً اعلیٰ نوع دی
منطق دی نحو نوں شروع
توں سیدھا متعین کرنا
زیادہ صاف ہوندا اے۔
مثال لئی، محدود تشکیل
والیاں انواع دا مجموعہ
استقرائی طور تے ایس
طرح متعین کر سکدے آں:
\begin{enumerate}
\item $\Nat$ اک محدود تشکیل والی نوع اے۔
\item جے $\sigma$ تے $\tau$ محدود تشکیل
  والیاں انواع ہون، تاں $\sigma
  \to \tau$ وی محدود تشکیل والی نوع اے۔
\item جے $\sigma$ تے $\tau$ محدود تشکیل
  والیاں انواع ہون، تاں $\sigma \times
  \tau$ وی محدود تشکیل والی نوع اے۔
\end{enumerate}
ایتھے محدود ہون دا مطلب نوع دی تشکیل دے محدود
قدم نیں؛ نوع دیاں شےواں دی گنتی متناہی ہون دا
دعویٰ نہیں۔
مطلب ایہ اے کہ $\Nat$
قدرتی اعداد دی نوع،
$\sigma \to \tau$ نوع
$\sigma$ توں نوع $\tau$
ول فنکشناں دی نوع،
تے $\sigma \times \tau$
جوڑیاں دی نوع اے
جیہناں دا اک رکن
$\sigma$ توں تے دوجا
$\tau$ توں ہوئے۔
فیر اصطلاحاں دا
مجموعہ استقرائی طور
تے ایس طرح متعین
کر سکدے آں:
\begin{enumerate}
\item ہر نوع $\sigma$ لئی
  متغیراں $x$, $y$,
  $z$, \dots دا
  ذخیرہ موجود اے؛
  ایہ نوع $\sigma$
  دے متغیر نیں۔
\item $\Obj 0$ نوع $\Nat$
  دی اصطلاح اے۔
\item $\Obj S$ (جانشین)
  نوع $\Nat \to \Nat$
  دی اصطلاح اے۔
\item جے $s$ نوع $\sigma$
  دی اصطلاح، تے $t$
  نوع $\Nat \to (\sigma \to \sigma)$
  دی اصطلاح ہوئے، تاں
  $\Obj{R}_{st}$ نوع
  $\Nat \to \sigma$
  دی اصطلاح اے۔
\item جے $s$ نوع
  $\tau \to \sigma$ دی
  اصطلاح تے $t$
  نوع~$\tau$ دی اصطلاح
  ہوئے، تاں $s(t)$
  نوع $\sigma$ دی
  اصطلاح اے۔
\item جے $s$ نوع~$\sigma$
  دی اصطلاح تے $x$
  نوع~$\tau$ دا متغیر
  ہوئے، تاں $\lambd[x][s]$
  نوع $\tau \to \sigma$
  دی اصطلاح اے۔
\item جے $s$ نوع~$\sigma$
  دی اصطلاح تے $t$
  نوع~$\tau$ دی اصطلاح
  ہوئے، تاں $\tuple{s, t}$
  نوع $\sigma \times
  \tau$ دی اصطلاح اے۔
\item جے $s$ نوع~$\sigma \times \tau$
  دی اصطلاح ہوئے، تاں
  $p_1(s)$ نوع~$\sigma$
  دی تے $p_2(s)$
  نوع~$\tau$ دی اصطلاح
  اے۔
\end{enumerate}
مطلب ایہ اے کہ
$\Obj{R}_{st}$ ایس
بازگشتی تعریف والے
فنکشن نوں بیان کردا اے:
\begin{align*}
\Obj{R}_{st}(0) & = s \\
\Obj{R}_{st}(x+1) & = t(x, R_{st}(x)),
\end{align*}
% PNB-SEM-032: arrow-type دے مطابق recursor دی دو دلیل والی لکھت دو لگاتار اطلاقاں دا مختصر نشان اے۔
بازگشت دی دوجی مساوات وچ دو دلیل والی لکھت
دو لگاتار اطلاقاں دا مختصر نشان اے: پہلے قدرتی عدد
اُتے اطلاق نال اک فنکشن ملدا اے، فیر اوہ پچھلی
حاصل قدر اُتے لاگو ہوندا اے۔
$\tuple{s, t}$ ایس جوڑے
نوں بیان کردا اے جیہدا
پہلا رکن~$s$ تے دوجا
رکن~$t$ اے؛ تے $p_1(s)$
تے~$p_2(s)$،~$s$ دے
پہلے تے دوجے رکن
(”پروجیکشن“) نوں بیان
کردے نیں۔ آخرکار
$\lambd[x][s]$ اوہ
فنکشن~$f$ بیان کردا
اے جیہڑا ایویں
متعین اے:
\[
f(x) = s
\]
ہر~$x$ لئی، جیہڑا
% PNB-SYN-011: شق6 وچ lambda دا input tau تے body/result sigma اے؛ اصل دی وضاحت وچ غلط sigma نوں tau کیتا۔
نوع~$\tau$ دا ہوئے؛
ایس طرح شق (6)
خاصیت توں فنکشن بنان
دا اک روپ دیندی اے،
جیہدے نال اصطلاحاں
راہی فنکشن متعین
ہو سکدے نیں۔
!!^{formula}s اکّو
نوع دیاں اصطلاحاں
وچکار !!{identity}
دے بیان $\eq[s][t]$،
عام قضیاتی رابطیاں
تے اعلیٰ نوع دے
مسوّراں توں بنائے
جاندے نیں۔ فیر
نظام دے مسلّمے
اوپر متعین اصطلاحاں
دیاں بنیادی مساوات،
تے نال مسوّراں
تے !!{identity}
دے عام منطقی
قاعدے رکھ سکدے
آں۔

جے محدود تشکیل والیاں انواع دے
نظام وچ سچائی دیاں
قدراں دی نوع
$\Omega$ شامل کر
دتی جاوے، تاں
ایہدی ورتوں نوں
چلاؤن والے مسلّمے
وی شامل کرنے
پین گے۔ اصل وچ
ذرا ہنر نال
پیچیدہ !!{formula}s
نوں بالکل ہٹا
کے انہاں دی
تھاں نوع $\Omega$
دیاں اصطلاحاں
رکھیاں جا
سکدیاں نیں!{}
ثبوتی نظام
فیر اوس حساب
نال بدل
سکدا اے۔
نتیجہ بنیادی
طور تے
\emph{انواع دا
سادہ نظریہ}
اے جیہڑا
الونزو چرچ
نے 1930 دی
دہائی وچ
پیش کیتا۔

دوجے درجے
دی منطق
وانگ، اعلیٰ
نوع دی معنیات
دے وی وکھ
وکھ روپ
ورتے جا
سکدے نیں۔
مکمل روپ
وچ نوع
$\sigma \to \tau$
دے متغیر
نوع~$\sigma$
دیاں شےواں
توں نوع~$\tau$
دیاں شےواں
تک دے
\emph{سارے}
فنکشناں
اُتے چل
دے نیں۔
جداں امید
ہو سکدی
اے، ایہ
معنیات
اینی طاقتور
اے کہ
ایدے لئی
کوئی تمام،
مؤثر
!!{derivation}
نظام نہیں۔
پر اک
کمزور معنیات
وی وچار
سکدے آں،
جیہدے وچ
!!a{structure}
عناصر دے
سیٹ
$T_\tau$،
ہر نوع
$\tau$ لئی
اک، تے نال
اطلاق،
پروجیکشن
وغیرہ دے
مناسب عمل
رکھدی اے۔
جے تفصیل
ٹھیک طرح
پوری کیتی
جاوے، تاں
اوپر بیان
کیتے
!!{derivation}
نظاماں
واسطے
تمامیت دے
قضیے مل
سکدے نیں۔

اعلیٰ نوع
دی منطق
پرکشش اے
کیوں جے
ایہ ایسا
ڈھانچا دیندی
اے جیہدے
وچ ریاضی
دا وڈا حصہ
قدرتی طریقے
نال سمویا
جا سکدا اے:
$\Nat$
توں شروع
کرکے حقیقی
اعداد،
متصل فنکشن
وغیرہ متعین
ہو سکدے
نیں۔ ایہ
وجدانیت
پسند منطق
دے سیاق
وچ وی خاص
طور تے
پرکشش اے،
کیوں جے
انواع دیاں
”تعمیری“
تشریحاں
صاف نیں۔
بلکہ اعلیٰ
نوع دی
معنیات دے
تعمیری روپ
وی تیار
کیتے جا
سکدے نیں
(جیہڑے
کلاسیکی
منطق دی
بجائے
وجدانیت
پسند منطق
اُتے مبنی
ہون)؛
ایہ روپ
انہاں
تعمیری
تشریحاں
نوں چنگی
طرح واضح
کردے نیں
تے کئی
لحاظاں
نال کلاسیکی
ہم منصباں
نسبت ودھ
دلچسپ نیں۔

\end{document}
